\section{Introduction and scope}\label{sec:introduction}

This paper makes a deliberately limited prime-weighted transfer.  The input
from the preceding papers is an exact finite character expansion for a
centered Chinese-remainder pair kernel, including its conductor multipliers
and a Parseval identity after averaging fully active shifts
\citep{WangCenteredKernels2026,WangPowerShortRoughBoxes2026}.  The analytic
inputs used here are classical: the multiplicative large sieve and the
Siegel--Walfisz theorem.  We assemble these ingredients first for arbitrary
coefficients and then for von Mangoldt weights.  We do not prove a new large
sieve inequality, a new Siegel--Walfisz estimate, or a new theorem on primes
in arithmetic progressions.

Let
\begin{equation}\label{eq:intro-Q-G-kappa}
 Q=\prod_{w<p\le w^2}p,
 \qquad
 G=(\Z/Q\Z)^\times,
 \qquad
 g=\varphi(Q),
 \qquad
 \kappa=\prod_{p\mid Q}\frac{p-2}{p-1}.
\end{equation}
For intervals $I,J\subset\Z$ of respective lengths $M,N<Q$, coefficient
sequences $a=(a_m)_{m\in I}$ and $b=(b_n)_{n\in J}$, and $(h,Q)=1$, define
\begin{equation}\label{eq:intro-weighted-box}
 \cB_h(a,b)
 =\sum_{\substack{m\in I,\ n\in J\\(mn,Q)=1}}
 a_m b_n
 \left\{
  \kappa^{-1}\one_{(mn+h,Q)=1}-1
 \right\}.
\end{equation}
The centering in \eqref{eq:intro-weighted-box} is the exact complete-group
centering.  It is independent of the locations of the intervals and of the
chosen coefficients.  Put
\begin{equation}\label{eq:intro-coefficient-energies}
 E_a=\sum_{\substack{m\in I\\(m,Q)=1}}|a_m|^2,
 \qquad
 E_b=\sum_{\substack{n\in J\\(n,Q)=1}}|b_n|^2.
\end{equation}
For $\chi\in\widehat G$, write $q_\chi$ for its primitive conductor and
\begin{equation}\label{eq:intro-weighted-transforms}
 A_{I,a}(\chi)
 =\sum_{\substack{m\in I\\(m,Q)=1}}a_m\chi(m),
 \qquad
 A_{J,b}(\chi)
 =\sum_{\substack{n\in J\\(n,Q)=1}}b_n\chi(n).
\end{equation}
The exact kernel coefficient attached to a nonprincipal character has
modulus
\begin{equation}\label{eq:intro-conductor-multiplier}
 d(\chi)=\prod_{p\mid q_\chi}\frac1{p-2}
 \ll \frac1{q_\chi}.
\end{equation}
This multiplier, rather than a generic estimate for a periodic kernel, is
what permits conductor-by-conductor aggregation below.

\subsection{Coefficient-uniform transfer}

The first layer does not use the arithmetic nature of the coefficients.  If
$|a_m|,|b_n|\le1$, exact-conductor orthogonality and the fully active shift
Parseval identity give
\begin{equation}\label{eq:intro-bounded-complete-shift}
 \frac1{M^2N^2}
 \frac1g\sum_{\ell\in G}|\cB_{2\ell}(a,b)|^2
 \ll
 \frac1{\max\{M,N\}}+\frac1{w\log w}.
\end{equation}
The estimate is uniform in the intervals and in all bounded coefficient
sequences.  It is an average over shifts $2\ell$ for which every prime
dividing $Q$ remains active; it is not a conclusion for each $\ell$.
The exact identity and its coefficient-uniform consequence are stated in
\cref{prop:weighted-shift-parseval,thm:weighted-complete-shift}.

The complete shift mean can again be sampled on a subpower interval.  If
$\cL$ is any interval of
\begin{equation}\label{eq:intro-subpower-shift-length}
 H=\left\lfloor Q^{1/\log w}\right\rfloor
\end{equation}
consecutive integers and
$H_Q=\#\{\ell\in\cL:(\ell,Q)=1\}$, the deterministic completion argument
transfers \eqref{eq:intro-bounded-complete-shift} to the average over those
$H_Q$ shifts.  Although $H=Q^{o(1)}$, this remains a mean-square statement.
In particular, it does not select a prescribed member of $\cL$.
See \cref{cor:weighted-short-shift}.

At a fixed shift, the character phases no longer diagonalize.  A dyadic
conductor decomposition and the classical multiplicative large sieve yield
the separate estimate
\begin{equation}\label{eq:intro-weighted-fixed-shift}
 |\cB_h(a,b)|
 \ll
 \sqrt{E_aE_b}
 \left\{
  \frac{\sqrt{MN}}w
  +(\log Q)(\sqrt M+\sqrt N)
  +\sqrt Q
 \right\},
 \qquad (h,Q)=1.
\end{equation}
The bound is uniform in the unit shift $h$.  Consequently, bounded
coefficients have normalized fixed-shift cancellation whenever $M,N$ are
fixed positive powers of $Q$ and $MN/Q\to\infty$.  A convenient uniform
form is $M,N\ge Q^\delta$ and $MN\ge Q^{1+\delta}$ for fixed $\delta>0$.
The ambient-volume threshold $MN=Q$ is the range furnished by this argument,
not an asserted barrier.  The precise uniform estimate is
\cref{thm:weighted-fixed-shift-large-sieve}.

The large-sieve mechanism in \eqref{eq:intro-weighted-fixed-shift} is worth
making explicit.  On a dyadic conductor block $T<q_\chi\le2T$, the factor
$d(\chi)\ll T^{-1}$ and the multiplicative large sieve control the weighted
second moment of each transform.  Cauchy--Schwarz is then applied within the
same conductor block.  Summing the resulting geometric blocks up to a
cutoff and using full character Parseval above that cutoff gives
\eqref{eq:intro-weighted-fixed-shift}; choosing the cutoff near $\sqrt Q$
balances the two sides.  Thus the estimate is not a pointwise bound for an
individual prime-supported character sum.

\subsection{Double von Mangoldt weights}

We next take dyadic factor intervals
\begin{equation}\label{eq:intro-dyadic-prime-boxes}
 I_X=(X,2X]\cap\Z,
 \qquad
 I_Y=(Y,2Y]\cap\Z,
 \qquad X,Y<Q,
\end{equation}
and insert $a_m=\Lambda(m)$ and $b_n=\Lambda(n)$.  Write the resulting form
as $\cB_h^{\Lambda,\Lambda}(X,Y)$.  If
\begin{equation}\label{eq:intro-power-sized-factors}
 X=Q^{\lambda+o(1)},
 \qquad
 Y=Q^{\nu+o(1)}
\end{equation}
for fixed $0<\lambda,\nu<1$, then
\begin{equation}\label{eq:intro-lambda-complete-rms}
 \frac1{XY}
 \left(
  \frac1g\sum_{\ell\in G}
  |\cB_{2\ell}^{\Lambda,\Lambda}(X,Y)|^2
 \right)^{1/2}
 =o(1).
\end{equation}
The same asymptotic conclusion holds on every shift interval of length
\eqref{eq:intro-subpower-shift-length}.  Hence all but $o(H_Q)$ of its unit
shift parameters satisfy the corresponding $o(XY)$ estimate.  This
density-one conclusion still does not identify $\ell=1$ or any other
prescribed shift.  These statements are
\cref{thm:lambda-complete-shift-rms,cor:lambda-short-shift-rms,cor:lambda-density-one-shifts}.

The fixed-shift large-sieve estimate supplies a complementary result.  For
every $(h,Q)=1$, including $h=2$, we obtain
\begin{equation}\label{eq:intro-lambda-fixed-result}
 \cB_h^{\Lambda,\Lambda}(X,Y)=o(XY)
 \qquad\text{when}\qquad \lambda+\nu>1.
\end{equation}
Uniform versions are stated in terms of fixed lower power bounds for $X,Y$
and a fixed power separation of $XY$ from $Q$.  The distinction between
\eqref{eq:intro-lambda-complete-rms} and
\eqref{eq:intro-lambda-fixed-result} is structural: shift orthogonality
allows every pair of positive power lengths in the mean-square theorem,
whereas the present fixed-shift large-sieve argument requires $XY$ to exceed
$Q$ by a power; see \cref{thm:lambda-fixed-shift}.

At small conductor, the prime-weighted input is the classical
Siegel--Walfisz theorem.  Uniformly for nonprincipal primitive characters of
conductor at most a fixed power of $\log X$, it gives arbitrary logarithmic
savings for the twisted summatory von Mangoldt function; see
\citet{Siegel1935,Walfisz1936} and the modern formulation in
\citet[Corollary~5.29]{IwaniecKowalski2004}.  Restricting to $(n,Q)=1$ does
not create a new prime-character-sum problem here.  Since $\Lambda$ is
supported on prime powers, the discrepancy from the unrestricted primitive
twist is bounded by
\begin{equation}\label{eq:intro-rough-lambda-comparison}
 \sum_{p\mid Q}\sum_{\substack{k\ge1\\X<p^k\le2X}}\log p
 \le \sum_{p\mid Q}\log p=\log Q,
\end{equation}
because $p>2$ makes a dyadic interval contain at most one power of each
$p$.  This is negligible on the power scales in
\eqref{eq:intro-power-sized-factors}.  At complete and short shift means,
larger conductors are controlled by the summable exact multiplier tail.  At
a prescribed shift, they are aggregated by the multiplicative large sieve of
\citet{MontgomeryVaughan1973}; standard accounts include
\citet[Chapter~7]{IwaniecKowalski2004}.  The elementary bounds
$\sum_{n\le x}\Lambda(n)\ll x$ and
$\sum_{n\le x}\Lambda(n)^2\ll x\log x$ provide the required coefficient
mass and energy.  None of these analytic estimates is new here.

Removing prime powers from \eqref{eq:intro-lambda-fixed-result} gives a
concrete prime-supported consequence.  In its fixed-shift range,
\begin{equation}\label{eq:intro-prime-prime-rough-asymptotic}
 \sum_{\substack{X<p\le2X,\ Y<q\le2Y\\(pq+h,Q)=1}}
 (\log p)(\log q)
 =(\kappa+o(1))XY.
\end{equation}
Uniform partial summation gives the corresponding unweighted asymptotic
with main term
$\kappa XY/(\log X\log Y)$.  The condition on the target is exactly the one
displayed: $pq+h$ avoids the primes in the finite window defining $Q$.
The weighted and unweighted forms are recorded in
\cref{cor:prime-prime-rough-target}.

\subsection{Relation to classical prime-distribution methods}

The multiplicative large sieve is a classical mean-square inequality for
arbitrary coefficients; see the foundational developments
\citep{Bombieri1965,DavenportHalberstam1966,Gallagher1967,MontgomeryVaughan1973}.
Our use of it is specialized only through the exact multiplier
\eqref{eq:intro-conductor-multiplier} and the fact that exact local support
equals primitive conductor for the squarefree modulus $Q$.  We make no
claim that the large-sieve inequality itself, or its application to a bare
von Mangoldt transform, is new.

This is also not a new Bombieri--Vinogradov theorem.  The classical theorem
controls prime-distribution errors on average over a range of moduli
\citep{Bombieri1965,Vinogradov1965}; variants for other multiplicative
sequences require their own hypotheses and exceptional-character analysis
\citep{GranvilleShao2018}.  Here $Q$ is one prescribed squarefree product.
In \eqref{eq:intro-lambda-complete-rms} the average is over local shifts, not
over moduli, while \eqref{eq:intro-lambda-fixed-result} is obtained from a
weighted conductor decomposition internal to that single $Q$.  These are
different averaging geometries.

Vaughan's identity decomposes $\Lambda$ into Type~I and Type~II pieces
\citep{Vaughan1977}, but the identity alone estimates neither piece.  The
present proofs do not establish a new prime-sensitive Type~II or dispersion
estimate.  Such estimates would become necessary if one replaced the local
roughness indicator by a prime detector at the target, or tried to continue
past the fixed-shift ranges supplied by \eqref{eq:intro-weighted-fixed-shift}.
Likewise, results on prime-supported character sums in short progressions,
such as \citet{Puchta2003}, address a different averaging and parameter
regime and are not silently imported here.

The large-modulus theorems of
\citet{BombieriFriedlanderIwaniec1986} and the fixed-residue and
well-factorable results of
\citet{Maynard2025FixedResidue,Maynard2025WellFactorable} use additional
structure in the residue class, modulus family, or modulus weights.  That
structure is not supplied merely by factoring $Q$ into primes or by observing
that $d(q)$ is multiplicative.  Those theorems are context for a possible
future dispersion layer, not inputs that automatically strengthen the
claims made here.

\subsection{Strict boundary of the conclusions}

The two factor variables in \eqref{eq:intro-prime-prime-rough-asymptotic}
are prime, but the target $pq+h$ is only $Q$-rough.  In particular, even at
$h=2$ the theorem does not assert that $pq+2$ is prime.  It is not an
estimate for
\[
 \sum_{n\le x}\Lambda(n)\Lambda(n+2),
\]
and it gives no lower bound for twin primes.  The shift-mean statements
cannot be specialized to $h=2$ without the independent fixed-shift theorem,
and the latter still concerns a semiprime-type factor product followed by a
finite roughness test.

Finally, local divisibility information does not remove the parity
obstruction of sieve theory; see \citet{HeathBrown1982} and the broader sieve
framework in \citet{FriedlanderIwaniec2010}.  Prime-producing asymptotic
sieves require further bilinear information absent here
\citep{FriedlanderIwaniec1998}.  Accordingly, we do not describe
\eqref{eq:intro-prime-prime-rough-asymptotic} as evidence for the twin-prime
conjecture and make no priority claim.  The contribution is narrower: an
exact CRT kernel converts classical Siegel--Walfisz and multiplicative
large-sieve input into coefficient-uniform fixed-shift bounds and
prime-weighted power-short shift means, with every remaining prime-detection
interface left explicit.
